\documentclass[UTF8,12pt,a4paper,fontset=fandol]{ctexart}

% 支持从项目根目录或当前笔记目录编译。
\makeatletter
\def\input@path{{../}{../../}}
\makeatother
\usepackage{researchnotes}

\title{三角行列式的性质与证明}
\author{}
\date{}

\begin{document}
\maketitle

\section{问题与性质}

一般的行列式由许多带符号的乘积项组成，但当主对角线的一侧全为零时，
其值只由主对角线元素决定。原因在于：\impo{每个乘积项都必须从每行、每列各取一个元素，}
而三角结构使主对角线以外的所有选取方式都含有零因子。
下面先从行列式的排列定义证明这一点，再给出按行、列展开的归纳证明。

全文设 $n$ 为正整数，$\mathbb F=\mathbb R$ 或 $\mathbb C$，
$A=(a_{ij})\in\mathbb F^{n\times n}$，其中 $a_{ij}$ 表示第 $i$ 行、第 $j$ 列的元素。
矩阵 $A$ 的\keyterm{行列式}（determinant）记为 $\det A$，也记为 $|A|$；
这里的竖线表示行列式，不表示绝对值。

\begin{definition}[三角矩阵]
元素 $a_{11},a_{22},\ldots,a_{nn}$ 所在的位置构成
\keyterm{主对角线}（main diagonal）。
若 $i>j$ 时 $a_{ij}=0$，则称 $A$ 为\keyterm{上三角矩阵}（upper triangular matrix）；
若 $i<j$ 时 $a_{ij}=0$，则称 $A$ 为\keyterm{下三角矩阵}（lower triangular matrix）。
两者统称为\keyterm{三角矩阵}（triangular matrix），其行列式称为三角行列式。
\end{definition}

上、下三角矩阵分别具有下列形式：
\[
  \begin{pmatrix}
    a_{11}&a_{12}&\cdots&a_{1n}\\
    0&a_{22}&\cdots&a_{2n}\\
    \vdots&\ddots&\ddots&\vdots\\
    0&\cdots&0&a_{nn}
  \end{pmatrix},
  \qquad
  \begin{pmatrix}
    a_{11}&0&\cdots&0\\
    a_{21}&a_{22}&\ddots&\vdots\\
    \vdots&\vdots&\ddots&0\\
    a_{n1}&a_{n2}&\cdots&a_{nn}
  \end{pmatrix}.
\]
主对角线及另一侧的元素不受限制，也允许为零。

\begin{theorem}[三角行列式的值]
\label{thm:triangular-determinant}
若 $A$ 是上三角矩阵或下三角矩阵，则
\begin{equation}
  \det A=\prod_{i=1}^{n}a_{ii}=a_{11}a_{22}\cdots a_{nn}.
  \label{eq:triangular-product}
\end{equation}
该结论不要求主对角线元素非零。
\end{theorem}

\section{证明一：从排列定义出发}
\label{sec:permutation-proof}

\subsection{展开式中的每一项如何选取元素}

记 $S_n$ 为集合 $\{1,2,\ldots,n\}$ 的所有\keyterm{排列}（permutation）组成的集合。
一个排列 $\sigma$ 是该集合到自身的双射，因此
$\sigma(1),\ldots,\sigma(n)$ 恰好是 $1,\ldots,n$ 的一种排列。
其\keyterm{逆序数}（inversion number）定义为
\[
  \tau(\sigma)
  =\#\{(i,j):1\leq i<j\leq n,\ \sigma(i)>\sigma(j)\},
\]
其中 $\#$ 表示集合中元素的个数。由行列式的定义，
\begin{equation}
  \det A
  =\sum_{\sigma\in S_n}(-1)^{\tau(\sigma)}
    \prod_{i=1}^{n}a_{i,\sigma(i)}.
  \label{eq:leibniz-formula}
\end{equation}
乘积 $\prod_{i=1}^{n}a_{i,\sigma(i)}$ 从第 $i$ 行取第 $\sigma(i)$ 列的元素。
因为所有 $\sigma(i)$ 互不相同，所以它恰好从每行、每列各取一个元素。

对上三角矩阵，可以先从最后一行理解零因子的作用：
最后一行只有第 $n$ 列的元素可能非零，因此要使乘积非零，就必须选 $a_{nn}$。
第 $n$ 列被占用后，倒数第二行只能选 $a_{n-1,n-1}$；如此逐行向上，
只有主对角线的选法可能留下。下三角矩阵则可以从第一行开始作同样的判断。
下面用一个关于排列的引理把这一思路写成严格证明。

\subsection{排列的约束与完整证明}

\begin{lemma}
\label{lem:one-sided-permutation}
设 $\sigma\in S_n$。若对所有 $i$ 都有 $\sigma(i)\geq i$，
或者对所有 $i$ 都有 $\sigma(i)\leq i$，则 $\sigma(i)=i$ 对所有 $i$ 成立。
\end{lemma}

\begin{proof}
因为 $\sigma(1),\ldots,\sigma(n)$ 是 $1,\ldots,n$ 的重排，所以
\[
  \sum_{i=1}^{n}\bigl(\sigma(i)-i\bigr)
  =\sum_{i=1}^{n}\sigma(i)-\sum_{i=1}^{n}i=0.
\]
若各差值都是非负整数，只要其中一个为正，总和就为正，故它们必须全部为零。
若各差值都是非正整数，只要其中一个为负，总和就为负，故也必须全部为零。
两种情形均得 $\sigma(i)=i$。
\end{proof}

\begin{proof}[定理~\ref{thm:triangular-determinant} 的排列证明]
\textbf{上三角情形。}
设 $A$ 为上三角矩阵。对于任意不是恒等排列的 $\sigma$，
由引理~\ref{lem:one-sided-permutation}，不可能对所有 $i$ 都有 $\sigma(i)\geq i$。
因此至少存在一个 $i$ 使 $\sigma(i)<i$。
此时 $a_{i,\sigma(i)}$ 位于主对角线下方，等于 $0$，从而
\[
  \prod_{i=1}^{n}a_{i,\sigma(i)}=0.
\]
所以式~\eqref{eq:leibniz-formula} 中，除恒等排列外的各项全部为零。
恒等排列 $\operatorname{id}$ 满足 $\operatorname{id}(i)=i$，其逆序数为 $0$，
故该项的符号为 $(-1)^0=1$，于是
\[
  \det A=a_{11}a_{22}\cdots a_{nn}.
\]

\textbf{下三角情形。}
设 $A$ 为下三角矩阵。对于任意不是恒等排列的 $\sigma$，
引理~\ref{lem:one-sided-permutation} 保证至少存在一个 $i$ 使 $\sigma(i)>i$。
此时 $a_{i,\sigma(i)}$ 位于主对角线上方，等于 $0$，所以对应的乘积项也为零。
展开式仍然只剩下恒等排列对应的项，得到同一公式。
\end{proof}

这里“只剩下”是指其他项都已证明为零，并不表示主对角线项一定非零。
若某个 $a_{ii}=0$，那么最后保留的这一项也等于 $0$，证明仍然有效。

\section{证明二：按行、列展开与数学归纳法}

若已经学过行列式的按行、列展开公式，还可以通过降阶证明该性质。
当 $n\geq2$ 时，记 $M_{ij}$ 为从 $A$ 中删去第 $i$ 行和第 $j$ 列后所得的
$(n-1)$ 阶行列式，称为元素 $a_{ij}$ 的\keyterm{余子式}（minor）。
按第一列、第一行展开的公式分别为
\begin{align}
  \det A&=\sum_{i=1}^{n}(-1)^{i+1}a_{i1}M_{i1},
  \label{eq:first-column-expansion}\\
  \det A&=\sum_{j=1}^{n}(-1)^{1+j}a_{1j}M_{1j}.
  \label{eq:first-row-expansion}
\end{align}
本节以这两个展开公式为已知结论；第~\ref{sec:permutation-proof} 节的证明不依赖它们。

\begin{proof}[定理~\ref{thm:triangular-determinant} 的归纳证明]
对矩阵的阶数 $n$ 作\keyterm{数学归纳法}（mathematical induction）。
当 $n=1$ 时，$\det(a_{11})=a_{11}$，结论成立。

设 $n\geq2$，并假设结论对所有 $n-1$ 阶上三角矩阵和下三角矩阵均成立。
令 $B$ 为删去 $A$ 的第一行、第一列后所得的矩阵，因此 $M_{11}=\det B$。

若 $A$ 为上三角矩阵，则第一列除 $a_{11}$ 外全为零。
由式~\eqref{eq:first-column-expansion}，
\[
  \det A=(-1)^{1+1}a_{11}M_{11}=a_{11}\det B.
\]
此时 $B$ 仍为上三角矩阵，主对角线元素依次为 $a_{22},\ldots,a_{nn}$，
故由归纳假设得 $\det B=a_{22}\cdots a_{nn}$，从而得到式~\eqref{eq:triangular-product}。

若 $A$ 为下三角矩阵，则第一行除 $a_{11}$ 外全为零。
由式~\eqref{eq:first-row-expansion}，同样有 $\det A=a_{11}\det B$。
此时 $B$ 仍为下三角矩阵，归纳假设仍给出 $\det B=a_{22}\cdots a_{nn}$。
因此下三角情形的结论也成立，归纳完成。
\end{proof}

这一证明每次只是把 $a_{11}$ 乘在余子式前面，没有除以 $a_{11}$，
因此即使 $a_{11}=0$，降阶等式仍成立，无须另作非零假设。

\section{例题与使用要点}

\begin{example}[非对角元素不影响三角行列式的值]
对任意 $x,y,z\in\mathbb F$，计算
\[
  D=\begin{vmatrix}
    2&x&y\\
    0&-3&z\\
    0&0&4
  \end{vmatrix}.
\]
\end{example}

\begin{solve}
这是上三角行列式，由式~\eqref{eq:triangular-product}，
$D=2\cdot(-3)\cdot4=-24$。
只要保持主对角线下方全为零，改变 $x,y,z$ 就不会改变行列式的值。
\end{solve}

\begin{example}[对角元素可以为零]
计算
\[
  D=\begin{vmatrix}
    1&0&0\\
    2&0&0\\
    3&4&5
  \end{vmatrix}.
\]
\end{example}

\begin{solve}
这是下三角行列式，故 $D=1\cdot0\cdot5=0$。
虽然没有一整行或一整列全为零，行列式仍然可以为零。
\end{solve}

对三角矩阵，由式~\eqref{eq:triangular-product} 及实数、复数乘积的性质可知：
行列式非零当且仅当每个主对角线元素都非零。
使用这一公式的前提是先确认矩阵具有三角结构；对一般矩阵，不能只乘主对角线元素。
例如
\[
  \begin{vmatrix}1&2\\3&4\end{vmatrix}
  =1\cdot4-2\cdot3=-2\ne1\cdot4.
\]
其中非主对角线项没有零因子，因而不能从展开式中略去。

\end{document}
